\newpage
\section{Сертификационные центры США}

Американская система оценки соответствия робототехнических систем строится на взаимодействии государственных исследовательских институтов, разрабатывающих методики испытаний, и независимых организаций, выполняющих оценку соответствия готовых изделий. Характерной чертой американского подхода является ориентация на сквозные системные стандарты, которые позволяют оценивать безопасность робота как готового изделия, снижая необходимость раздельного тестирования его отдельных компонентов.

Ведущую роль в разработке единых измерительных метрик играет государственная лаборатория Национального института стандартов и технологий (NIST Robotics Test Facility), расположенная в Гейтерсберге. Площадь внутренних помещений центра составляет 892 квадратных метра. Совместно с профильными комитетами специалисты института разрабатывают стандартизированные методы испытаний мобильности, манипуляций и человеко-машинного взаимодействия, охватывающие более 50 сценариев для наземных, воздушных и водных систем (\url{https://www.nist.gov/laboratories/tools-instruments/robotics-test-facility}). Важно, что данные испытания ориентированы на проверку робота в сборе как единого комплекса, выполняющего целевую задачу.

В коммерческом секторе сертификацию по принципу единого окна осуществляет корпорация Intertek. Вместо проведения десятков разрозненных тестов на механику, электропроводку и программное обеспечение по отдельности, центр проводит испытания на соответствие сквозному стандарту ANSI/CAN/UL 3100. Этот документ оценивает безопасность робота в целом: функциональную безопасность систем управления, надежность бортовых зарядных устройств, защиту двигателей от перегрузки и взрывобезопасность литиевых батарей в рамках единого испытательного цикла (\url{https://www.intertek.com/industrial/robotics/}).

Глобальный сертификационный центр UL Solutions реализует аналогичный комплексный подход на базе стандарта UL 4600, регламентирующего безопасность автономных продуктов. Вместо проверки соответствия разрозненным техническим требованиям к отдельным деталям, стандарт требует от разработчика подготовить единое доказательство безопасности системы. Данный документ подробно обосновывает надежность всей системы в сборе при ее эксплуатации в заданной среде. Дополнительно UL Solutions проводит сертификацию систем искусственного интеллекта по программе UL 3115, оценивающей устойчивость и безопасность алгоритмов в составе готового изделия (\url{https://www.ul.com/services/certification}).

Ассоциация развития автоматизации дополняет систему путем оценки готовых робототехнических ячеек и комплексов в сборе на соответствие стандарту ANSI/A3 R15.06-2025. Этот стандарт объединяет требования к манипуляторам, защитным ограждениям и оптическим датчикам, позволяя оценивать безопасность промышленной роботизированной ячейки как единого комплекса (\url{https://www.automate.org/robotics/robotics-certifications/robotic-integrator-certification-program}).

\begin{longtblr}[
    caption = {Сравнительный анализ инфраструктуры сертификационных центров США},
    label = {tab:usa_centers},
]{
    colspec = { X[1,l,m] X[1,l,m] X[2,l,m] X[2,l,m] },
    hlines, vlines,
    rowhead = 1,
    rows = {font=\small},
    row{1} = {font=\bfseries\small},
}
    Организация & Системные стандарты & Доступная инфраструктура и методы & Ограничения / Дефициты \\
    NIST Robotics Test Facility & ASTM E54.08 & Изолированные камеры (температура, свет), полигоны с сыпучими грунтами, испытательные стенды для гуманоидных роботов & Не предоставляет услуги коммерческой сертификации серийной продукции \\
    Intertek & ANSI/CAN/UL 3100 & Испытания бортовых зарядных станций, камеры электромагнитной совместимости, стенды проверки литиевых батарей & Зависимость от методик сторонних разработчиков, отсутствие крупных открытых полигонов \\
    UL Solutions & UL 4600, UL 3115 & Электротехнические лаборатории, программные комплексы аудита кибербезопасности нейросетевых моделей & Смещение фокуса с механических испытаний на аудит доказательной документации \\
    Association for Advancing Automation & ANSI/A3 R15.06-2025 & Выездной аудит готовых ячеек, проверка квалификации интеграторов & Отсутствие собственной физической испытательной базы \\
\end{longtblr}

\newpage
\section{Сертификационные центры Европейского союза}

Европейская инфраструктура сертификации ориентирована на выдачу единого сквозного подтверждения безопасности готового изделия — маркировки CE. Вместо прохождения независимых проверок у разных ведомств, производитель сертифицирует робота в рамках единого регламента машинной безопасности, привлекая для системной оценки нотифицированные органы.

Крупнейшим нотифицированным органом в Европейском союзе является немецкий концерн TÜV SÜD. Центр проводит комплексную сертификацию беспилотного складского транспорта и мобильных роботов по сквозному международному стандарту ISO 3691-4. Данный стандарт оценивает безопасность робота в сборе: логику систем управления, динамическое торможение, работу сканеров безопасности и поведение машины при потере связи, избавляя от необходимости раздельно сертифицировать корпус, датчики и приводные системы (\url{https://www.tuvsud.com/en/industries/manufacturing/robotics-and-mobile-robotics}).

Прикладную валидацию комплексных систем осуществляет Испытательная лаборатория робототехники при институте Fraunhofer IPA в Штутгарте. Лаборатория специализируется на проверке сервисных и медицинских роботов в сборе по стандарту ISO 13482. Тестирование проводится на передовых платформах в условиях, полностью имитирующих бытовые помещения и клиники, что позволяет оценить безопасность траекторий движения робота и его манипуляторов в непосредственном контакте с человеком (\url{https://www.ipa.fraunhofer.de/en/about-us/laboratories_and_testing_facilities/Robotics_Testing_Facility.html}).

Международная корпорация SGS со штаб-квартирой в Швейцарии обеспечивает системную оценку робототехнической продукции по Директиве по радиооборудованию. Экспертиза SGS охватывает проверку стабильности беспроводных каналов управления, помехозащищенности и общей электробезопасности готового изделия, а также проведение комплексных испытаний аккумуляторных систем по стандарту ООН UN/DOT 38.3 (\url{https://www.sgs.com/en/services/robotic-products-testing-and-certification}).

\begin{longtblr}[
    caption = {Сравнительный анализ инфраструктуры сертификационных центров Европейского союза},
    label = {tab:eu_centers},
]{
    colspec = { X[1,l,m] X[1,l,m] X[2,l,m] X[2,l,m] },
    hlines, vlines,
    rowhead = 1,
    rows = {font=\small},
    row{1} = {font=\bfseries\small},
}
    Организация & Системные стандарты & Доступная инфраструктура и методы & Ограничения / Дефициты \\
    TÜV SÜD & ISO 3691-4, ISO 10218-1/2 & Программно-аппаратные стенды симуляции отказов, зоны тестирования безопасного взаимодействия, камеры электромагнитной совместимости & Высокая стоимость услуг, длительные сроки проведения комплексных испытаний \\
    Fraunhofer IPA & ISO 13482 & Специализированные платформы, передовые испытательные и демонстрационные стенды для практических испытаний в лаборатории & Ориентация на научно-исследовательские задачи в ущерб потоковой коммерческой сертификации \\
    SGS & Machinery Regulation, RED, UN/DOT 38.3 & Стенды разрушающего контроля батарей, программные комплексы тестирования уязвимостей беспроводных сетей & Ограниченные возможности по тестированию крупногабаритной робототехники на открытых полигонах \\
\end{longtblr}

\newpage
\section{Сертификационные центры Китая}

Китайская модель регулирования робототехники ориентирована на исключение избыточных процедур за счет внедрения национальной системы сертификации CR. Этот сквозной знак качества подтверждает безопасность и надежность всего робототехнического комплекса в целом, заменяя для производителя необходимость прохождения десятков ведомственных проверок.

Базовую инфраструктуру для выдачи сертификатов CR формирует государственная сеть национальных центров тестирования и оценки роботов, распределенная по ключевым промышленным кластерам страны. Крупнейший центр в Шанхае оснащен климатическими и вибрационными стендами для испытаний промышленных манипуляторов и тяжелых беспилотных тележек в сборе на соответствие единым национальным стандартам качества КНР (\url{https://www.baclcorp.com.cn/show.asp?para=en_6_994_713}).

Для системной валидации ИИ-компонентов в Пекине создан специализированный инновационный кластер, включающий Пекинский инновационный центр робототехники воплощённого интеллекта. Вместо раздельной проверки программного обеспечения и датчиков, центр оценивает работоспособность систем машинного зрения и алгоритмов принятия решений непосредственно в составе готовых роботов, проходящих испытания в зоне экономического развития Пекина (\url{https://english.wuhan.gov.cn/H_1/NWP/202602/t20260211_2728662.shtml}).

Инновационный центр гуманоидных роботов провинции Хубэй в Ухане площадью 7000 квадратных метров ориентирован на оценку сложных человекоподобных платформ. Вместо разрозненных тестов, робот оценивается интегрально в рамках 23 прикладных сценариев (бытовых и производственных). База данных центра генерирует до 24 000 единиц данных ежедневно, что позволяет проверить стабильность ИИ-модели управления роботом во всех штатных режимах (\url{https://www.yicaiglobal.com/news/chinas-wuhan-unveils-usd139-million-fund-to-lead-in-humanoid-robots}, \url{http://www.factoryautmation.com/xinwefile/Robot-authentication-is-coming_v001232.html}).

\begin{longtblr}[
    caption = {Сравнительный анализ инфраструктуры сертификационных и инновационных центров Китая},
    label = {tab:china_centers},
]{
    colspec = { X[1,l,m] X[1,l,m] X[2,l,m] X[2,l,m] },
    hlines, vlines,
    rowhead = 1,
    rows = {font=\small},
    row{1} = {font=\bfseries\small},
}
    Организация & Системные стандарты & Доступная инфраструктура и методы & Ограничения / дефициты \\
    Национальные центры тестирования и оценки роботов & Схемы CR-сертификации & Стенды комплексных механических и электромагнитных испытаний промышленных манипуляторов и тяжелых платформ в сборе & Процедуры ориентированы прежде всего на внутренний рынок и национальные требования КНР \\
    Инновационный центр робототехники воплощённого интеллекта (Пекин) & Стандарты ИИ-систем, допуск готовых платформ & Лаборатории, инфраструктура для исследований, испытаний и внедрения решений в зоне экономического развития & Больше ориентирован на разработку и пилотирование, чем на финальную сертификацию \\
    Инновационный центр гуманоидных роботов (Хубэй) & Сквозное тестирование ИИ-моделей в сборе & Помещения площадью 7000 кв. м, 23 комплексных сценария применения готовых систем & Узкая специализация, не покрывает весь спектр испытаний промышленной робототехники \\
\end{longtblr}

\newpage
\section{Сертификационные центры Южной Кореи, Японии, Сингапура и ОАЭ}

В странах Азии и на Ближнем Востоке допуск робототехники строится на базе комплексных полигонных тестов готовых изделий и сертификации по сквозным международным стандартам безопасности сервисных машин.

Японский центр безопасности роботов (в составе JARI) проводит системную оценку мобильных сервисных машин и экзоскелетов по единому международному стандарту ISO 13482. Прохождение этого стандарта гарантирует безопасность изделия в целом за счет тестов на устойчивость, предотвращение столкновений и ограничение силы при физическом контакте с человеком. Ассоциация качества Японии (JQA) стала первой организацией в мире, выдавшей комплексный сертификат соответствия по стандарту ISO 13482:2014. Для испытаний беспилотного транспорта используется полигон Shirosato Test Center, где оценивается поведение систем помощи водителю на скользких трассах и в искусственной городской застройке (JARI, \url{https://www.jari.or.jp/}).

В Сингапуре центр передового опыта CETRAN реализует сквозной допуск беспилотного транспорта по интегральной методике Milestone 1 и Milestone 2. Вместо аудита отдельных плат, датчиков и программных модулей, беспилотная платформа тестируется как единый комплекс в условиях имитационного города. Успешное выполнение тестов (включая проезд перекрестков, обход препятствий и движение под проливным дождем) позволяет получить автоматическое разрешение на допуск машины на дороги общего пользования (CETRAN, \url{https://cetran.sg/}).

Южнокорейские центры, включая Корейскую испытательную лабораторию, используют международную систему взаимного признания CB Scheme для трансграничного признания единого сертификата соответствия без повторных испытаний. Основное внимание лаборатория уделяет проверке безопасности взаимодействия человека и мобильного робота-курьера на базе стандартов ISO, используя городские испытательные зоны и государственные программы регуляторных песочниц в Сеуле.

Испытательные полигоны Управления по дорогам и транспорту в ОАЭ (Дубай) ориентированы на комплексную проверку беспилотных такси и логистических платформ в сборе на устойчивость к экстремальным температурам (выше 45--50 °C) и запыленности. Оценивается способность готовой системы компенсировать деградацию сенсоров при песчаных бурях без прерывания миссии.

\begin{longtblr}[
    caption = {Сравнительный анализ инфраструктуры сертификационных центров Азии и Ближнего Востока},
    label = {tab:asia_centers},
]{
    colspec = { X[1,l,m] X[1,l,m] X[2,l,m] X[2,l,m] },
    hlines, vlines,
    rowhead = 1,
    rows = {font=\small},
    row{1} = {font=\bfseries\small},
}
    Организация & Системные стандарты & Доступная инфраструктура и методы & Ограничения / Дефициты \\
    JARI / JQA (Япония) & ISO 13482, ISO 10218-1/2 & Салазки для краш-тестов, антропоморфные манекены, полигон Shirosato Test Center, климатические и экранированные комнаты & Ограниченное покрытие задач тяжелой промышленной робототехники \\
    CETRAN (Сингапур) & Протокол Milestone 1/2 & Многоуровневые дорожные полигоны, системы имитации тропических ливней, манекены пешеходов & Фокус преимущественно на городской мобильности, ограниченная площадь полигонов \\
    Испытательные центры Южной Кореи & Международная система CB Scheme, ISO 13482 & Испытательные зоны взаимодействия человека и робота, городские тестовые полигоны, регуляторные песочницы & Зависимость от временных разрешений в рамках экспериментальных правовых режимов \\
    Управление по дорогам и транспорту (ОАЭ) & Системный климатический допуск RTA & Термокамеры, песчаные треки, испытательные стенды устойчивости готовых сенсорных систем беспилотных такси & Основной акцент на климатических испытаниях, меньший охват сценариев сложного взаимодействия с пешеходами \\
\end{longtblr}

\newpage
\section{Сравнительный анализ международных и отечественных центров}

Отечественные центры сертификации и испытаний (ФГУП «НАМИ», ПромМаш Тест, ВНИИС, Университет Иннополис) обладают развитой материальной базой для подтверждения соответствия робототехники требованиям базовых Технических регламентов Таможенного союза. Однако при переходе к комплексной сертификации готовых робототехнических изделий возникает существенный технологический разрыв. В Российской Федерации отсутствуют национальные стандарты, эквивалентные сквозным международным стандартам ISO 3691-4 (для безэкипажного складского транспорта) или ISO 13482 (для сервисной мобильной робототехники), способным подтвердить безопасность робота как единого изделия в рамках одной процедуры.

В результате отечественные разработчики вынуждены проходить разрозненные проверки отдельных узлов машины на разных площадках. В то же время зарубежные центры предлагают комплексные полигонные тесты и системный аудит алгоритмов автономного вождения в рамках единой методологии. Сравнительные характеристики ведущих центров представлены в \cref{tab:certification_centers_comparison}.

\begin{longtblr}[
    caption = {Сравнительный анализ международных и отечественных сертификационных центров},
    label = {tab:certification_centers_comparison},
]{
    colspec = { X[1,l,m] X[1,l,m] X[2,l,m] X[2,l,m] },
    hlines, vlines,
    rowhead = 1,
    rows = {font=\small},
    row{1} = {font=\bfseries\small},
}
    Организация (Страна) & Системные стандарты & Доступная инфраструктура & Ограничения / Дефициты \\
    TÜV SÜD (Германия) & ISO 3691-4, ISO 10218-1/2, Machinery Regulation & Программно-аппаратные стенды симуляции отказов, зоны тестирования безопасного взаимодействия, камеры электромагнитной совместимости & Высокая стоимость услуг, жесткая ориентация строго на стандарты и директивы Европейского союза \\
    JARI / JQA (Япония) & ISO 13482 & Краш-тесты, антропоморфные манекены, крупный полигон с имитацией города & Узкая специализация исключительно на автомобильных и сервисных мобильных платформах \\
    CETRAN (Сингапур) & Протокол Milestone 1/2 & Имитационный город, многоуровневые полигоны, дождевые установки & Ограниченность рамками городского автономного транспорта, отсутствие промышленного сегмента \\
    ФГУП «НАМИ» (РФ) & ТР ТС 018/2011, ТР ТС 010/2011 & Дорожные полигоны, климатические камеры, динамические тесты беспилотных транспортных средств & Отсутствие методик и специализированных стендов для комплексных испытаний коботов (1*) \\
    ПромМаш Тест (РФ) & ТР ТС 010/2011, ТР ТС 020/2011 & Сертификационные стенды, механические и электромагнитные испытания & Отсутствие крупных открытых полигонов и тестов высокоуровневой автономности (1*) \\
    Университет Иннополис (РФ) & Внутренние методики тестирования навигации & Тестовые ячейки, системы навигации, симуляционные среды & Отсутствие аккредитации для проведения обязательной государственной сертификации \\
\end{longtblr}
{\small\noindent
    \textit{Примечание: (1*) — отсутствие оборудования для оценки функциональной безопасности высокого уровня, кибербезопасности и сложных сценариев взаимодействия готовых робототехнических изделий.}
}

Анализ показывает, что ни один отечественный центр на данный момент не способен предоставить полный комплекс услуг по системной валидации искусственного интеллекта и машинного зрения, что затягивает вывод отечественной продукции на коммерческий рынок.

\newpage
\section{Возможности применения зарубежного опыта при создании ИЦСРТ}

Проведенный анализ международного опыта демонстрирует, что для обеспечения технологического суверенитета и ускоренного вывода на рынок безопасной робототехники необходимо концептуальное изменение подходов к оценке соответствия. Формируемый Испытательный центр сертификации робототехнических технологий (ИЦСРТ) должен отказаться от фрагментированного тестирования компонентов и стать единым методологическим хабом, внедряющим комплексную сертификацию робота в сборе.

Ключевой задачей ИЦСРТ должно стать внедрение сквозных методик испытаний, аналогичных китайской системе CR-сертификации или американским стандартам семейств UL 3100 и UL 4600. Это позволит отечественным производителям подтверждать безопасность складских, промышленных и сервисных роботов в рамках одного комплексного испытательного цикла. Элементы международного опыта, рекомендуемые к поэтапному внедрению в структуру ИЦСРТ, представлены в \cref{tab:icst_implementation}.

\begin{longtblr}[
    caption = {Элементы зарубежного опыта для внедрения в ИЦСРТ},
    label = {tab:icst_implementation},
]{
    colspec = { X[2,l,m] X[1,l,m] X[2,l,m] X[2,l,m] X[2,l,m] },
    hlines, vlines,
    rowhead = 1,
    rows = {font=\small},
    row{1} = {font=\bfseries\small},
}
    Практика & Страна-донор & Возможность применения в ИЦСРТ & Ограничения & Ожидаемый эффект \\
    Единая сквозная сертификация робототехнических систем (по типу китайской CR-сертификации) & Китай & Разработка национальных регламентов комплексной оценки безопасности робота в сборе & Требует изменения нормативной базы на уровне ЕАЭС & Сокращение количества необходимых сертификатов с десятка мелких до одного единого знака качества \\
    Стенды для оценки биомеханического воздействия (по стандарту ISO/TS 15066) & Япония, ЕС & Закупка специализированных манекенов и датчиков давления для проверки безопасного контакта коботов с человеком & Отсутствие утвержденных национальных ГОСТов для проведения измерений & Легализация и масштабирование безопасного внедрения коботов на отечественных заводах \\
    Виртуальная валидация готовых изделий и полунатурное моделирование & США, ЕС & Интеграция цифровых двойников для предварительной проверки навигации и логики работы робота до вывода на полигон & Сложность точной валидации физических моделей лидарных и ультразвуковых сенсоров & Снижение сроков и стоимости испытаний комплексных систем в сборе на 30--40\% \\
    Системная сертификация автономных платформ (по типу стандартов UL 3100 и UL 4600) & США & Разработка сквозного стандарта безопасности для складских автономных платформ, объединяющего тесты электроники, ПО и механики & Противоречит жесткой парадигме ведомственной оценки & Появление правовой возможности легализации роботов с ИИ и машинным обучением \\
\end{longtblr}

Создание ИЦСРТ с применением указанных подходов позволит ликвидировать фрагментацию отечественной испытательной базы, обеспечив полный цикл системной сертификации робототехнических комплексов в сборе в режиме единого окна.